% Источники: exam/materials/преподаватель/2020/23-03-2020-MODEL L5(2).doc; exam/materials/моделирование.pdf, раздел 32.
\section{Критерий согласия Колмогорова–Смирнова.}

\subsection*{Постановка задачи}

Проверяется нулевая гипотеза $H_0$: выборка $x_1,\ldots,x_n$ взята из генеральной совокупности с теоретической функцией распределения $F_0(x)$.

\subsection*{Алгоритм}

\begin{enumerate}
  \item \textbf{Построить эмпирическую функцию распределения} (ECDF):
    \[
      F_n(x) = \frac{1}{n}\sum_{i=1}^n \mathbf{1}[x_i \le x].
    \]
    Практически: упорядочить выборку $x_{(1)} \le x_{(2)} \le \cdots \le x_{(n)}$; $F_n(x_{(k)}) = k/n$.

  \item \textbf{Вычислить наблюдаемую статистику}:
    \[
      D_n = \sup_{x}\bigl|F_n(x) - F_0(x)\bigr|.
    \]
    Для упорядоченной выборки удобно считать максимум по скачкам ECDF:
    \[
      D_n^+ = \max_{1\le k\le n}\left(\frac{k}{n}-F_0(x_{(k)})\right), \qquad
      D_n^- = \max_{1\le k\le n}\left(F_0(x_{(k)})-\frac{k-1}{n}\right),
    \]
    \[
      D_n = \max(D_n^+,D_n^-).
    \]

  \item \textbf{Найти критическое значение} $D_n^*(\beta)$ при уровне значимости $\beta$:
    \begin{itemize}
      \item при $n > 35$: \quad $D_n^*(\beta) \approx \sqrt{-\tfrac{1}{2}\ln\beta} \cdot \tfrac{1}{\sqrt{n}}$;
      \item при $n \le 35$: значение берётся из таблицы критических значений критерия Колмогорова.
    \end{itemize}

  \item \textbf{Принятие решения}:
    \begin{itemize}
      \item если $D_n < D_n^*(\beta)$ — $H_0$ \emph{не отвергается};
      \item если $D_n \ge D_n^*(\beta)$ — $H_0$ \emph{отвергается}.
    \end{itemize}
\end{enumerate}

\subsection*{Пример}

Пусть $n = 5$, $\beta = 0.05$, проверяется нормальность.
Выборка: $x = \{1.1,\; 2.0,\; 2.5,\; 3.1,\; 4.2\}$.
Оценки параметров: $\hat{\mu}=2.58$, $\hat{\sigma}=1.08$.

\medskip
\noindent
\begin{tabular}{cccc}
  \hline
  $k$ & $x_{(k)}$ & $F_n(x_{(k)})$ & $F_0(x_{(k)})$ \\
  \hline
  1 & 1.1 & 0.20 & 0.22 \\
  2 & 2.0 & 0.40 & 0.43 \\
  3 & 2.5 & 0.60 & 0.56 \\
  4 & 3.1 & 0.80 & 0.72 \\
  5 & 4.2 & 1.00 & 0.91 \\
  \hline
\end{tabular}

\medskip
$D_n^+ = 0{,}09$, $D_n^- = 0{,}22$, поэтому $D_n = 0{,}22$.

По таблице ($n=5$, $\beta=0.05$): $D_5^* = 0.565$.
Так как $D_n = 0{,}22 < 0{,}565$, $H_0$ \emph{не отвергается} — выборка согласуется с нормальным распределением.

\clearpage
